\def\stat{l1-9}





\def\tit{СРЕДСТВА ПОДДЕРЖКИ ИСПОЛНЯЕМОГО КОДА, СИНТЕЗИРОВАННОГО 


ПО~СПЕЦИФИКАЦИЯМ, НА~ЯЗЫКЕ~CELL}





\def\titkol{Средства поддержки исполняемого кода на языке Cell}


% синтезированного  по~спецификациям на языке Cell} 





\def\autkol{О.\,А. Бондаренко, К.\,И. Волович, В.\,А. Кондрашев}





\def\aut{О.\,А. Бондаренко$^1$, К.\,И. Волович$^2$, В.\,А. Кондрашев$^3$}





\titel{\tit}{\aut}{\autkol}{\titkol}





%{\renewcommand{\thefootnote}{\fnsymbol{footnote}}\footnotetext[1]


%{Работа выполнена при частичной финансовой поддержке по Программе фундаментальных 


%исследований ОНИТ РАН на 2011~г.\ (проект~1.5).}}





\renewcommand{\thefootnote}{\arabic{footnote}}


\footnotetext[1]{Институт проблем информатики Российской академии наук, olga@ipi.ac.ru}


\footnotetext[2]{Институт проблем информатики Российской академии наук, kv@ipi.ac.ru}


\footnotetext[3]{Институт проблем информатики Российской академии наук, vd@ipi.ac.ru}





     


       \Abst{В статье рассматриваются аспекты генерации 


служебного программного кода, обеспечивающего функционирование 


протокольных автоматов, разработанных на языке Cell. 


Описывается понятие <<исполняющей сис\-те\-мы>> как совокупности 


программного кода, включаемого компилятором языка Cell в 


процедурный код с целью обеспечения функционирования основного 


алгоритма протокольного автомата. Сериализация процедурного 


кода, выполняемая компилятором Cell во время его синтеза, 


предполагает возможность эффективного функционирования 


автомата без использования операционной сис\-те\-мы (ОС), что 


накладывает на исполняющую систему задачи управления всеми 


ресурсами, используемыми автоматом.}


       


       \KW{язык Cell; ячейка; исполняющая сис\-те\-ма; управление 


ресурсами; сериализация; иерархический автомат; синтез 


программного обеспечения; телекоммуникационный протокол}





  \vskip 14pt plus 9pt minus 6pt








      \thispagestyle{headings}





      


            \label{st\stat}


            


            


            


\vspace*{-12pt}


      


       


\section{Введение}


      


      Современные телекоммуникационные сис\-те\-мы функционируют на 


базе широкого спектра протоколов, каждый из которых описывается 


соответствующим стандартом. Как правило, описание 


телекоммуникационного протокола в таком стандарте состоит из 


формальной (блок-схе\-мы алгоритмов, описание блоков данных) и 


неформальной час\-ти, содержащей словесное описание функционирования 


сис\-те\-мы. Создание на основе такого описания программного обеспечения, 


реализующего телекоммуникационный протокол, является творческой 


задачей для ин\-же\-не\-ра-про\-грам\-миста, фактически не поддающейся 


автоматизации. 


      


      Язык Cell, предлагаемый авторами в качестве средства разработки 


про\-грам\-мно\-го обеспечения телекоммуникационных протоколов, дает 


возможность описать алгоритм функционирования компонентов 


телекоммуникационной сис\-те\-мы формализованным способом.  Наличие 


такого языка позволяет производить автоматизированный синтез 


программного обеспечения путем трансляции формальных спецификаций 


на языке Cell на язык реализации, подходящий для использования на 


конкретных вычислительных платформах. В~качестве языка реализации 


авторами выбран язык~C. 


      


      Однако простой трансляции кода протокольного автомата на язык 


реализации, а затем в машинные коды недостаточно для построения 


функционирующей сис\-те\-мы. Необходимо также решить вопросы 


управления распределением памяти, времени, других ресурсов 


вычислительной платформы, межпроцессного взаимодействия.


      


      Для решения данных вопросов обычно применяются два подхода:


      \begin{enumerate}[(1)]


      \item исполнение кода с использованием ОС 


как средства управления ресурсами и предоставления сервиса 


приложениям;


      \item исполнение целевого кода непосредственно вычислительной 


платформой без использования ОС.


      \end{enumerate}


      


      В первом случае все функции по управлению памятью, 


распределением процессорного времени между процессами, 


межпроцессному взаимодействию могут быть возложены на службы 


ОС. Во втором случае эти функции должны быть 


включены в код протокольного автомата. 


      


      Каждый из подходов имеет свои достоинства и недостатки в 


условиях применения во встраиваемых сис\-те\-мах. Использование 


ОС исключает необходимость подключения 


дополнительного кода на этапе компиляции автомата, что облегчает 


разработку компилятора языка Cell. Однако операционная сис\-те\-ма требует 


дополнительных ресурсов со стороны оборудования, что в условиях 


встраиваемых систем является существенным. Кроме того, механизм 


обеспечения параллельности (псевдопараллельности) функционирования 


компонентов протокольного автомата средствами ОС (процессы, нити) не 


требуется в рамках концепции сериализации, заложенной в основу языка 


Cell. Более того, для обеспечения механизма последовательного 


исполнения кода, сгенерированного по спецификациям языка Cell, в 


условиях функционирования под управлением ОС потребуется 


использование механизма блокировок, семафоров и иных средств 


ОС, обеспечивающих параллельность исполнения,


т.\,е.\ использование именно тех средств организации вычислительного 


процесса, от которых позволяет отказаться концепция сериализации.


      


      Второй подход, предусматривающий включение кода исполняющей 


сис\-те\-мы непосредственно в код протокольного автомата, позволяет 


полностью реализовать механизмы сериализации. Управление 


одновременностью исполнения логически независимых частей 


протокольного автомата в этом случае обеспечивается самим кодом 


протокольного автомата. Параллельно функционирующие логические 


модули в данном случае не оформляются как  отдельные объекты в рамках 


ОС (процессы, нити), а выделяются исключительно в 


рамках логики исполнения одного процесса протокольного автомата, 


исполняемого последовательно на аппаратных средствах встраиваемой 


сис\-те\-мы.


      


      Исходя из вышеизложенных посылок, авторы предлагают при 


генерации целевого кода протокольного автомата использовать второй 


подход, т.\,е.\ включать в состав протокольного автомата элементы 


исполняющей сис\-те\-мы, обес\-пе\-чи\-ва\-ющих функционирование без 


использования~ОС. 


      


      Данный подход обеспечивает мобильность кода между различными 


аппаратными платформами, поскольку не содержит аппаратно зависимых 


компонентов. Привязка к конкретному оборудованию может 


осуществляться путем применения аппаратного драйвера, 


обеспечивающего прием и отправку низкозкоуровневых данных через 


буфер в памяти с использованием прерываний. Интерфейс взаимодействия 


с драйвером может предоставляться внешней по отношению к автомату 


библиотекой и не являться непосредственно принадлежностью 


протокольного автомата.


      


      Разработка алгоритмов функционирования исполняющей сис\-те\-мы 


является отдельной задачей при создании компилятора. Программный код 


исполняющей сис\-те\-мы, включаемый в результирующий код 


протокольного автомата, является частью библиотек компилятора. 


Взаимодействие собственно автоматного кода и библиотек обеспечивается 


через программный интерфейс, вызовы которого формируются на этапе 


компиляции.


      


      Далее в статье показаны структурные модули исполняющей 


сис\-те\-мы и алгоритмы их функционирования, обеспечивающие работу 


протокольного автомата без использования механизмов~ОС.


      


\section{Структура исполняющей сис\-те\-мы}


      


      Основной функцией исполняющей сис\-те\-мы является распределение 


ресурсов и управление исполнением независимыми частями автомата.


      


      Авторами предлагается алгоритм генерации целевого кода 


протокольного автомата, при котором компилятор языка Cell производит 


статическое выделение ресурсов для каждого объекта независимого 


исполнения (контекста процесса, очередей, области данных). Таким 


образом, первичное распределение ресурсов уже проведено на этапе 


периода компиляции. На этапе исполнения от диспетчера ресурсов 


исполняющей сис\-те\-мы требуется лишь активация того или иного 


статического объекта в памяти.


      


      Также к функциям исполняющей сис\-те\-мы относится управление 


очередями сигналов, обеспечивающее последовательное (сериализованное) 


исполнение процедурного кода. Можно выделить основные модули, 


составляющие исполняющую систему:


      \begin{enumerate}[(1)]


      \item модуль управления очередями;


      \item модуль управления контекстами;


      \item модуль управления реакциями автомата (модуль переходов);


      \item модуль управления областями данных сигналов;


      \item модуль контроля использования ресурсов;


      \item модуль обработки исключительных ситуаций.


      \end{enumerate}


      


\section{Алгоритмы функционирования исполняющей 


сис\-те\-мы}


      


      \subsection{Модуль управления очередями}


      


      Модуль управления очередями предоставляет автомату интерфейс к 


функциям помещения сигналов в очередь и извлечение сигналов из 


очередей. Основной программный цикл (message loop), интегрируемый 


компилятором в исходный код автомата, использует функции GetSignal и 


DispatchSignal для получения сигнала из внешней очереди и обработки его 


кодом протокольного автомата соответственно.


      


      Спецификой построения логики работы автомата является наличие 


двух очередей для внешних и внутренних сигналов. Внешние сигналы 


принимаются от ап\-па\-рат\-но-за\-ви\-си\-мо\-го драйвера и размещаются 


в очереди внешних сигналов, в то время  как внутренние сигналы 


генерируются самим автоматом и размещаются, соответственно, в очереди 


внутренних сигналов.


      


      Указанная пара очередей привязывается к контексту ячейки. Как 


показано в~[1], каждой корневой ячейке сопоставляется свой контекст. 


      


      При генерации программного кода компилятор интегрирует 


неявным образом код управления очередями в протокольный автомат, 


причем функции управления очередями в качестве одного из аргументов 


получают указатель на контекст, с которым ассоциируется корневая ячейка. 


Таким образом обеспечивается реентерабельность кода управления 


очередями и возможность работы с очередями как базовой ячейки, так и 


клонов.


      


      Область хранения сигналов (тело очередей) размещается в каждом 


контексте и представляет собой с точки зрения программирования массив 


структур, включающих в себя указатель на дескриптор сигнала, и поля 


аргументов, передаваемых вместе с сигналом. Функции управления 


очередью производят размещение сигналов в ту или иную часть очереди в 


зависимости от их типов~[2].


      


      Функция размещения сообщения в очереди копирует аргументы из 


области данных аппаратного драйвера в область данных автомата, 


обеспечивая тем самым отсутствие необходимости хранения данных в 


области аппаратных очередей.


      


      Функция извлечения сообщений из очереди производит 


копирование данных в выделенную область памяти контекста с 


последующим освобождением очереди, пересортировкой оставшихся 


сигналов, пересчетом индексов. 


      


      Глубина очередей не является конечной и не определяется на этапе 


компиляции. Выделение памяти для элементов очереди происходит 


динамически. По этой причине имеется необходимость включения в 


исполняемый код блока динамического управления распределением 


памяти. Блок динамического распределения памяти включается в модуль 


контроля управления ресурсами и в случае необходимости позволяет 


выделить память для объектов очереди из нераспределенной области 


(кучи), выделяемой на этапе компиляции.


      


      Очевидно, что определение размера кучи на этапе компиляции 


представляет собой трудноразрешимую задачу. Явных указаний в тексте 


протокольного автомата на возможную длину очередей не содержится, 


поэтому решение этой задачи носит эвристический характер. 


      


      Управление размером кучи в период компиляции производится  


специальными директивами компилятору, на основании которых он 


осуществляет статическое распределение памяти в сегменте данных. 


Проверку правильности распределения необходимо проводить путем 


верификации и отладки результирующего кода.


      


      Отметим, что контроль за переполнением очередей в период 


исполнения производится модулем контроля использования ресурсов, 


функции которого будут рассмотрены ниже. В~случае проведения 


верификации модуль контроля можно использовать  для мониторинга 


заполнения очередей каждого контекста.





%\vspace*{-3pt}


      


      \subsection{Модуль управления контекстами}


      


 %     \vspace*{-1pt}


      


      Модуль управления контекстами обеспечивает хранение полного 


списка контекстов корневых ячеек протокольного автомата. Список 


контекстов является статическим и определяется на этапе компиляции. 


      


      Основной структурной единицей исполняемого кода является 


контекст корневой ячейки. Исполняющая сис\-те\-ма содержит в своем 


составе диспетчер управ\-ле\-ния контекстами, обеспечивающий 


переключение между независимо функционирующими частями автомата.   


      


      Для  простоты будем считать, что контекст базовой ячейки всегда 


один, а множественность контекстов появляется при клонировании. 


      


      С целью сохранения ресурсов целевого вычислителя компилятор 


языка Cell генерирует для автоматов-клонов реентерабельный 


процедурный код, позволяющий использовать один экземпляр 


исполняемого кода для всех однотипных ав\-то\-ма\-тов-клонов.


      


      При этом исполнение кода того или иного клона обеспечивается 


переключением контекста средствами рассматриваемого модуля.


      


      Модуль управления контекстами представляет собой библиотеку 


периода исполнения, подключаемую к целевому коду протокольного 


автомата на этапе компиляции. В задачи модуля входит:


      \begin{itemize}


      \item  поддержание списка активных контекстов на каждый момент 


времени функционирования протокольного автомата;


      \item переключение контекста текущего процесса в зависимости от 


сигнала, обрабатываемого функцией DispatchSignal модуля управления 


очередями.


      \end{itemize}


      


      Для реализации указанных функций модуль управления 


контекстами должен хранить таблицу ассоциаций областей данных 


контекстов и идентификаторов процессов клонов. 


      


      При выборе из очереди внешних сигналов сигнала, адресованного 


конкретному клону, модуль управления контекстами настраивает скрытые 


переменные среды (генерируемые компилятором неявно) на конкретную 


область памяти в массиве контекстов. Дальнейшие вызовы функций-обработчиков 


производятся с передачей в них настроенного контекста и, 


соответственно, выполняются в рамках конкретного клона. Поскольку 


контексты содержат полный состав пользовательских данных локального и 


общего классов видимости, то  такое переключение позволяет функциям 


реализации переходов протокольного автомата работать на пространстве 


переменных конкретного клона до тех пор, пока очередной сигнал из 


очереди внешних сигналов не переключит контекст клона.


      


      Поскольку распределение памяти для хранения контекстов 


производится на этапе компиляции, то при управлении контекстами 


(активации и деактивации) каждой корневой ячейки выполняется проверка 


доступности ресурса и отсутствия конфликта между количеством 


запрашиваемых и выделенных при компиляции контекстов. Количество 


контекстов каждого типа настраивается с помощью директив компилятору 


и не может быть изменено в период исполнения. В~отличие от размера 


очередей сигналов количество клонов каждого типа можно оценить более 


точно, поскольку в структуре протокольного автомата каждый клон 


отвечает за некоторый сервис, а количество сервисов каждого типа можно 


определить на этапе проектирования автомата по его спецификациям.


%      


      Тем не менее, компилятор включает в протокольный автомат код 


проверки переполнения списка контекстов клонов. Контроль производится 


с помощью модуля контроля использования ресурсов. Аналогично случаю 


контроля заполнения очередей сигналов модуль контроля может 


использоваться для мониторинга утилизации памяти контекстов при 


верификации автомата. В~случае обнаружения переполнения модуль 


контроля генерирует исключительную ситуацию, которая также 


обрабатывается специализированной библиотекой, подключаемой к 


автомату на этапе компиляции.


{\looseness=-2





}


      


      Каждый контекст клона содержит собственные экземпляры 


очередей внешних и внутренних сигналов, поэтому функции отправки 


сигналов взаимодействуют с модулем управления контекстами с целью 


получения доступа  к нужной очереди. Взаимодействие осуществляется 


неявно за счет скрытого кода, интегрируемого на этапе компиляции, 


позволяющего транслировать идентификатор клона в адрес 


      струк\-ту\-ры-кон\-текста, содержащей конкретные очереди. 


Механизм установления указанной ассоциации скрывается от 


разработчика автомата на языке Cell, который в качестве 


      про\-цес\-са-по\-лу\-ча\-те\-ля сигнала использует идентификатор 


клона, возвращенный функцией клонирования. 


      


      \subsection{Модуль управления реакциями автомата (модуль 


переходов)}


      


      Модуль управления реакциями автомата обеспечивает поддержку 


актуального списка сигналов, обрабатываемых автоматом в текущий 


момент времени, а также вызов необходимой функции перехода по 


идентификатору сигнала.


      


      Список реакций протокольного автомата хранится в контексте в 


виде списка адресов функций переходов, возможных в текущий момент 


времени. Учитывая иерархичность протокольного автомата, на каждом 


уровне иерархии имеется свой набор возможных реакций. 


      


      Основная задача модуля~--- формирование нового списка активных 


сигналов в векторе перехода после выполнения каждой функции перехода. 


Для этого компилятор включает в тело каждой функции перехода код, 


обеспечивающий вызов функции, производящей модификацию вектора 


перехода. Данная функция удаляет из вектора реакции, определяющие 


предыдущее состояние автомата, и устанавливает текущие. 


      


      Модуль управления реакциями предоставляет модулю управления 


очередями интерфейс для обращения к функции перехода по 


идентификатору сигнала.


      


      Таким образом, модуль управления очередями выполняет алгоритм 


обработки сигнала путем выборки из очереди внешних сигналов, а затем 


обработки очереди внутренних сигналов до полного ее опустошения. При 


этом для каждого сигнала (внешнего или внутреннего) вызывается 


функция из состава модуля управления реакциями для поиска адреса 


перехода, а каждый переход завершается вызовом функции установки 


новых реакций.


      


      Указанный алгоритм функционирования модуля управления 


реакциями автомата позволяет постоянно поддерживать в контексте 


протокольного автомата полный набор реакций на всех уровнях иерархии. 


      


      При возникновении исключительных ситуаций модуль управления 


реакциями вызывает соответствующие функции модуля обработки 


исключительных ситуаций.


      


      Исключения периода исполнения могут возникать в следующих 


случаях:


      \begin{itemize}


      \item  отсутствие в векторе переходов функ\-ции-обра\-бот\-чи\-ка 


для данного состояния;


      \item попытка записи в вектор перехода нескольких 


      функ\-ций-обра\-бот\-чи\-ков для одного сигнала.


      \end{itemize}


      


      Вообще говоря, в нормальных условиях функционирования 


указанные ситуации не должны возникать  в период исполнения, 


поскольку на этапе компиляции производится проверка соответствия 


устанавливаемых реакций состоянию, в которое переходит протокольный 


автомат. Однако при нарушении логики функционирования 


протокольного автомата, описанного на языке Cell, вызванного 


зацикливанием, бесконечной рекурсией, другими ошибками периода 


исполнения, не диагностируемыми на этапе компиляции, такая ситуация 


возможна. Поэтому обработка исключительных ситуаций такого характера 


особенно важна в период отладки и верификации автоматного кода. 


      


      В режиме отладки функции модуля управления реакциями, а также 


функции обработки исключительных ситуаций могут быть 


переопределены  с целью получения диагностической информации. 


В~рабочем режиме  функционирования автомата исключительная 


ситуация, связанная со множественной записью вектора перехода, является 


критической и должны приводить к остановке или перезапуску корневой 


ячейки текущего автомата.


      


\subsection{Модуль управления областями данных сигналов}


      


      Модуль управления областями данных сигналов предназначен для 


хранения и обработки протокольных блоков данных (PDU~--- Protocol Data Unit), передаваемых 


внешними сигналами. 


      


      Основная задача модуля~--- предоставление доступа коду 


протокольного автомата к данным, содержащимся в PDU.


      


      Алгоритм обработки PDU предполагает размещение PDU в 


контексте протокольного автомата (базового или клона) и обеспечения к 


нему доступа процедурного кода из любой функции перехода.


      


      Код модуля обработки PDU не генерируется компилятором, а 


предоставляется разработчиком в качестве внешней программной 


библиотеки для компоновки с кодом протокольного автомата. В~состав 


библиотечных функций  входят функции для кодирования/\!раскодирования 


полей PDU, используемых протокольным автоматом. Протокольный 


автомат использует для доступа к данным интерфейс, предоставляемый 


библиотекой. В~функции раскодирования передается адрес области 


памяти, содержащей собственно PDU и адрес области для размещения 


декодированных данных. Структура области декодированных данных 


определяется в ходе описания переменных диспетчерского, общего или 


локального классов видимости в коде протокольного автомата~[2--4]. 


Таким образом, автомат получает возможность обращения к области 


декодированных данных как к структуре, определенной на этапе 


компиляции. 


      


      Для корректного обращения к областям декодированных данных 


разработчик должен следить, чтобы все обращения выполнялись после 


инициализации данных путем вызова функции раскодирования, причем 


вызов должен завершиться успешно.  


      


      Для размещения данных в существующий PDU используется 


функция кодирования, предоставляемая указанной библиотекой. Функция 


обеспечивает упаковку структурированных данных протокольного 


автомата в битовый поток и размещение его в указанной области памяти. 


Кодирование вызывается из тела перехода протокольного автомата, и 


обработка кодов завершения указанной функции возлагается на 


разработчика автомата на языке Cell. 


      


      Контроль использования ресурсов и отсутствие нарушений в 


использовании памяти возлагаются на саму библиотеку модуля 


управления областями данных PDU. С целью контроля границ областей 


памяти в функции библиотеки передается размер каждой области, 


подлежащей заполнению, вычисляемый на этапе компиляции. Функции 


модуля сами отслеживают выход за допустимые границы областей, в том 


числе и при рекурсивной обработке. В~случае недостаточности 


выделенной области функция раскодирования прекращает работы и 


возвращает код ошибки.


      


      При кодировании данных в существующий PDU производятся 


аналогичные проверки границ, и в случае некорректно выделенной области 


памяти возвращается ошибка периода исполнения. При этом сохранность 


предыдущих данных, расположенных в PDU, определяется реализацией 


функций библиотеки. В~общем случае можно считать, что старые данные, 


содержащиеся в PDU, при неуспешном завершении функции 


кодирования будут уничтожены.


      


      Таким образом, внешний по отношению к коду протокольного 


автомата модуль обеспечивает доступ кода автомата к распакованным 


структурированным данным, передаваемым с внешними сигналами, а 


также упаковку структурированных данных в PDU для передачи с 


использованием сигналов.


      


      \subsection{Модуль контроля использования ресурсов}


      


      Модуль контроля использования ресурсов предназначен для 


проверки границ областей памяти, распределенной на этапе компиляции. 


Как отмечалось выше, при функционировании протокольного автомата 


может возникнуть нехватка ресурсов, выделенных компилятором. 


      В~част\-ности, это относится к количеству контекстов той или иной 


корневой ячейки.


      


      Код модуля контроля использования ресурсов генерируется 


компилятором и включает в себя в качестве данных таблицы 


распределения ресурсов всех типов, распределенных компилятором. 


В~каж\-дую функцию, занимающую ресурс (например, клонирование) 


компилятор включает вызов функции проверки доступности 


запрашиваемого ресурса. В~случае наличия свободного ресурса модуль 


помечает его в таблице как занятый и возвращает код завершения, 


свидетельствующий о возможности использовать ресурс. В~случае 


отсутствия возможности выделить ресурс модуль возвращает код ошибки. 


Задача проверки кода ошибки ложится на разработчика автомата на языке 


Cell.


      


      Другой функцией модуля контроля ресурсов является динамическое 


выделение памяти для размещения сигналов в очередях. Выделение 


происходит неявным образом. В~каждую функцию отправки сигнала 


добавляется код вызова модуля с целью получения указателя на область 


данных для размещения сигнала. Модуль контроля выделяет из кучи 


запрашиваемое количество памяти, ассоциирует его в своих внутренних 


таблицах с текущим контекстом и возвращает указатель на область памяти. 


Таким образом, для каждого сигнала выделяется необходимое и 


достаточное количество памяти (с учетом передаваемых с сигналом 


данных). 


      


      Возвращение памяти в кучу производится при извлечении сигналов 


из очередей. В~функцию извлечения сигнала из очереди добавляется 


вызов функции, освобождающей память. Данная функция помечает блок 


данных как неиспользуемый и возвращает его в кучу для дальнейшего 


распределения.


      


      Необходимо отметить, что такой механизм распределения памяти 


неизбежно порождает фрагментацию. Это снижает эффективность 


функционирования и может приводить к невозможности выделения 


ресурсов для конкретного сигнала, несмотря на то, что общее количество 


памяти в куче достаточно, но разбито на большое количество мелких 


фрагментов, каждого из которых недостаточно для выделения 


достаточного непрерывного блока. 


      


      Для борьбы с фрагментацией традиционно применяются механизмы 


<<сборки мусора>> либо виртуальной памяти.


      


      Авторы считают, что программная реализация механизма 


виртуальной памяти, предусматривающая страничную организацию 


памяти, трансляцию виртуальных адресов в физические слишком 


громоздка для встраиваемых систем. Такая задача обычно выполняется 


средствами ОС, от которой авторы решили отказаться. 


Ряд процессоров предлагают аппаратную реализацию виртуальной памяти, 


однако использование такой возможности нарушило бы принцип 


мобильности и независимости кода от аппаратной платформы. 


      


      Реализация же сборки мусора не требует применения подобных 


решений и с точки зрения авторов более подходит для использования в 


коде протокольного автомата. Вместе с тем, необходимо отметить, что 


сборка мусора может потребовать больше ресурсов периода исполнения, 


чем виртуальная память.


      


      Программный код сборщика мусора включается компилятором в 


код автомата и вызывается каждый раз при извлечении сигналов из 


очереди. Однако производить сортировку кучи при каждом извлечении 


сигнала представляется нецелесообразным, поскольку приводит к 


большим затратам процессорного времени. 


      


      Для определения необходимости запуска сборки мусора 


используется алгоритм оценки фрагментации, позволяющий вычислить 


некоторый индекс (степень фрагментации памяти). В~случае превышения 


индексом порогового значения при очередном извлечении сигналов из 


очереди запускается механизм дефрагментации. Также механизм 


дефрагментации запускается в период простоя сис\-те\-мы.


      


      Чтобы иметь возможность осуществлять не только слияние 


соседних освобожденных участков кучи, но и перемещать используемые 


участки памяти, функции распределения памяти возвращают указатель 


не непосредственно на распределенную память, а на 


ячейку, содержащую указатель на распределенную память. Таким образом, 


сборщик мусора при сортировке имеет возможность переметить 


используемую область памяти, поменяв ее адрес в указанной ячейке. При 


этом не произойдет сбоя функционирования основного кода, поскольку все 


обращения к памяти производятся по указателям, а те, в свою очередь, 


поддерживаются в актуальном состоянии.


      


      Описанный механизм управления памятью скрыт от разработчика 


протокольного автомата на языке Cell и включается в код автомата 


компилятором.


      


      Отметим, что вычисление индекса фрагментации и его порогового 


значения носит эвристический характер. По умолчанию он принимается 


равным количеству фрагментов в куче, а порог срабатывания~--- половине 


общего количества возможных в системе сигналов. Очевидно, что такой 


алгоритм не является оптимальным, однако позволяет время от времени 


запускать дефрагментацию.


      


      Более оптимальное управление дефрагментацией достигается путем 


вы\-чис\-ле\-ния индекса с учетом количества и размера фрагментов и сильно 


зависит от самого протокольного автомата. Поэтому в исполняющей 


системе предусмот\-рен механизм переопределения функции подсчета 


индекса и установки порога срабатывания. Пользователь может 


определить свою функцию и включить ее в исполняемый код на этапе 


компоновки.


      


      В случае невозможности выделения памяти для размещения 


сигнала механизм дефрагментации вызывается принудительно, если же 


сигнал невозможно разместить и после дефрагментации, вызывается 


исключительная ситуация.


      


      \subsection{Модуль обработки исключительных ситуаций}


      


      Модуль обработки исключительных ситуаций вызывается из любой 


точки сис\-те\-мы в случае нештатного функционирования протокольного 


автомата.


      


      Как показано выше, исключительные ситуации могут возникать в 


случае нарушений в работе протокольного автомата или в случае нехватки 


ресурсов (в част\-ности, памяти для размещения сигналов в очередях).


      


      По умолчанию в случае возникновения исключительной ситуации 


модуль производит освобождение всех ресурсов, относящихся к контексту, 


вызвавшему исключение, и помечает данный контекст как свободный. 


Если указанные действия относятся к автомату-клону, то происходит 


рекурсивное освобождение всех ресурсов, относящихся к контекстам 


текущего и вышележащих клонов. Если исключительная ситуация 


возникает в контексте базовой ячейки, то происходят полное 


освобождение всех ресурсов и перезапуск автомата.


      


      Функции модуля обработки исключений могут быть 


переопределены пользователем для нужд отладки и верификации кода.


      


\section{Заключение}


      


      Поддержка исполняемого кода протокольного автомата, 


синтезированного по спецификациям, на языке Cell осуществляется путем 


интеграции компилятором специализированного скрытого кода в целевой 


код протокольного автомата. 


      


      Основные принципы функционирования указанного кода состоят в 


сокрытии механизмов управления ресурсами от разработчика программы 


на языке Cell и обеспечения функционирования протокольного автомата в 


условиях отсутствия~ОС. 


      


      Данное решение позволяет использовать протокольные автоматы, 


специфицированные на языке Cell в условиях встраиваемых систем, 


характеризующихся ограниченными ресурсами вычислительных платформ.





{\small\frenchspacing


{%\baselineskip=10.8pt


\addcontentsline{toc}{section}{Литература}


\begin{thebibliography}{9}


      


\bibitem{1-9}


\Au{Бондаренко Т.\,В., Волович К.\,И., Кондрашев~В.\,А.}


Язык Cell~--- инструмент для синтеза программного обеспечения 


многоуровневого телекоммуникационного протокола по частично 


формализованным спецификациям~/\!\!/ Вестник Воронежского 


государственного ун-та, 2011. Т.~7. №\,3. С.~120--125.





\bibitem{2-9}


\Au{Бондаренко Т.\,В., Бондаренко О.\,А., Волович~К.\,И., 


Кондрашев~В.\,А.} 


Сигнальный механизм языка Cell~/\!\!/ Системы и средства информатики. 


Вып.~20. №\,3.~--- М.: ИПИ РАН, 2010. С.~45--66.





\bibitem{3-9}


\Au{Бондаренко Т.\,В., Бондаренко О.\,А., Волович~К.\,И., 


Кондрашев~В.\,А.} 


Язык Cell: модель обработки клонов~/\!\!/ Системы и средства 


информатики.  Вып.~20. №\,3.~--- М.: ИПИ РАН, 2010. С.~67--81.





 \label{end\stat}





\bibitem{4-9}


\Au{Бондаренко Т.\,В., Бондаренко О.\,А., Волович~К.\,И., 


Кондрашев~В.\,А.} 


Базовая модель функционирования автомата в системе программирования 


Cell~/\!\!/ Системы и средства информатики. Вып.~20. №\,3.~--- М.: ИПИ РАН, 


2010. С.~82--97.





 \end{thebibliography}


}


}





      


%Executing system for code synthesized on the base of 


%specifications on the language Cell


      %


%Bondarenko O., Kondrashev V., Volovich K.





       %The article considers issues of the program code generation, 


%which provides the functioning of the hierarchical state machines 


%developed on the language Cell. The concept of ``executing system'' as a 


%program code included by compiler of the language Cell into procedural 


%code in order to provide the functioning of the main algorithm of the 


%hierarchical state machines is described. The serialization of the 


%procedural code which is executed by compiler Cell during its synthesis 


%offers an opportunity of the effective functioning of the state machine 


%without usage of operational system. That assigns a task on the executing 


%system to manage all the resources which uses the hierarchical state 


%machine.


        


%Keywords: language Cell, cell, executing system, resource management, 


%serialization, hierarchical state machine, software synthesis, 


%telecommunication protocol.


        


